\textbf{Phần 1. Trắc nghiệm lựa chọn}
\setcounterpageref{ex}{0}

\begin{ex} %[5D10N61]
Vận tốc đi bộ bình thường của một học sinh tiểu học khoảng
\choice
{\True $4\text{ km/h}$}
{$15\text{ km/h}$}
{$40\text{ km/h}$}
{$1\text{ km/h}$}
\loigiai{
Vận tốc đi bộ trung bình khoảng $4\text{ km/h}$.\\
\textbf{Đáp án của câu hỏi là A}
}
\end{ex}

\begin{ex} %[5D10N61]
Vận tốc trung bình của một chiếc xe buýt chạy trong nội thành thường khoảng
\choice
{$80\text{ km/h}$}
{$5\text{ km/h}$}
{\True $25\text{ km/h}$}
{$120\text{ km/h}$}
\loigiai{
Xe buýt trong nội thành có tốc độ trung bình khoảng $20 - 30\text{ km/h}$ do phải dừng đón trả khách.\\
\textbf{Đáp án của câu hỏi là C}
}
\end{ex}

\begin{ex} %[5D10H61]
Nhà bạn Nam cách trường $1{,}8\text{ km}$. Mỗi ngày Nam đi xe đạp với vận tốc $12\text{ km/h}$. Thời gian Nam đi từ nhà đến trường là
\choice
{$6$ phút}
{$12$ phút}
{$15$ phút}
{\True $9$ phút}
\loigiai{
$t = 1{,}8 : 12 = 0{,}15\text{ giờ} = 9$ phút.\\
\textbf{Đáp án của câu hỏi là D}
}
\end{ex}

\begin{ex} %[5D10H61]
Một gia đình dự định đi ô tô từ Hà Nội về quê cách đó $100\text{ km}$ với vận tốc ước tính $50\text{ km/h}$. Nếu muốn có mặt ở quê lúc $10$ giờ sáng thì gia đình phải xuất phát muộn nhất lúc mấy giờ?
\choice
{\True $8$ giờ}
{$7$ giờ}
{$8$ giờ $30$ phút}
{$9$ giờ}
\loigiai{
Thời gian đi: $100 : 50 = 2$ giờ.\\
Thời điểm xuất phát muộn nhất: $10 - 2 = 8$ giờ.\\
\textbf{Đáp án của câu hỏi là A}
}
\end{ex}

\begin{ex} %[5D10H61]
Bạn Mai tập chạy cự li $400\text{ m}$. Đồng hồ bấm giờ chỉ $1$ phút $20$ giây. Vận tốc chạy của Mai là
\choice
{$4\text{ m/s}$}
{$6\text{ m/s}$}
{\True $5\text{ m/s}$}
{$8\text{ m/s}$}
\loigiai{
$1$ phút $20$ giây $= 80$ giây.\\
$v = 400 : 80 = 5\text{ (m/s)}$.\\
\textbf{Đáp án của câu hỏi là C}
}
\end{ex}

\begin{ex} %[5D10H61]
Khoảng cách giữa hai cột mốc ven đường cao tốc là $100\text{ m}$. Một ô tô chạy qua hai cột mốc này mất đúng $4$ giây. Vận tốc của ô tô là
\choice
{$60\text{ km/h}$}
{\True $90\text{ km/h}$}
{$80\text{ km/h}$}
{$100\text{ km/h}$}
\loigiai{
Vận tốc: $100 : 4 = 25\text{ m/s} = 25 \times 3{,}6 = 90\text{ km/h}$.\\
\textbf{Đáp án của câu hỏi là B}
}
\end{ex}

\begin{ex} %[5D10V61]
Một người tập thể dục đi bộ một vòng quanh bờ hồ hình tròn có bán kính $200\text{ m}$ với vận tốc $4\text{ km/h}$. Thời gian người đó đi hết $1$ vòng (lấy $\pi = 3{,}14$) xấp xỉ là
\choice
{$12$ phút}
{$15$ phút}
{$24$ phút}
{\True $18{,}84$ phút}
\loigiai{
Chu vi bờ hồ: $2 \times 3{,}14 \times 200 = 1\,256\text{ (m)} = 1{,}256\text{ km}$.\\
Thời gian đi: $1{,}256 : 4 = 0{,}314\text{ giờ} = 18{,}84$ phút.\\
\textbf{Đáp án của câu hỏi là D}
}
\end{ex}

\begin{ex} %[5D10V61]
Hằng ngày bố chở An đi học bằng xe máy với vận tốc $30\text{ km/h}$ hết $16$ phút. Hôm nay do trời mưa, bố chỉ đi được với vận tốc $24\text{ km/h}$. Thời gian An đến trường hôm nay là
\choice
{$18$ phút}
{\True $20$ phút}
{$22$ phút}
{$25$ phút}
\loigiai{
Quãng đường từ nhà đến trường: $30 \times \dfrac{16}{60} = 8\text{ (km)}$.\\
Thời gian khi đi $24\text{ km/h}$: $8 : 24 = \dfrac{1}{3}\text{ giờ} = 20$ phút.\\
\textbf{Đáp án của câu hỏi là B}
}
\end{ex}

\begin{ex} %[5D10V61]
Một tàu cao tốc đi quãng đường $300\text{ km}$ mất $1$ giờ $40$ phút. Vận tốc trung bình của tàu cao tốc là
\choice
{\True $180\text{ km/h}$}
{$150\text{ km/h}$}
{$200\text{ km/h}$}
{$240\text{ km/h}$}
\loigiai{
$1$ giờ $40$ phút $= \dfrac{5}{3}$ giờ.\\
$v = 300 : \dfrac{5}{3} = 180\text{ (km/h)}$.\\
\textbf{Đáp án của câu hỏi là A}
}
\end{ex}

\begin{ex} %[5D10V61]
Một bạn học sinh ước lượng thời gian đi bộ từ tầng 1 lên tầng 4 của trường học (mỗi tầng cao $3{,}6\text{ m}$, có $20$ bậc thang mỗi tầng). Biết mỗi giây bạn bước được $2$ bậc thang. Thời gian bạn bước lên tầng 4 là
\choice
{$20$ giây}
{$40$ giây}
{$60$ giây}
{\True $30$ giây}
\loigiai{
Từ tầng 1 lên tầng 4 có $3$ nhịp cầu thang: $3 \times 20 = 60$ bậc thang.\\
Thời gian bước: $60 : 2 = 30$ giây.\\
\textbf{Đáp án của câu hỏi là D}
}
\end{ex}

\begin{ex} %[5D10V61]
Một đàn nhạn bay về phương nam tránh rét với vận tốc $45\text{ km/h}$. Mỗi ngày đàn nhạn bay liên tục trong $8$ giờ. Sau $5$ ngày đàn nhạn bay được quãng đường là
\choice
{$1\,600\text{ km}$}
{\True $1\,800\text{ km}$}
{$2\,000\text{ km}$}
{$2\,250\text{ km}$}
\loigiai{
Quãng đường mỗi ngày: $45 \times 8 = 360\text{ (km)}$.\\
Quãng đường 5 ngày: $360 \times 5 = 1\,800\text{ (km)}$.\\
\textbf{Đáp án của câu hỏi là B}
}
\end{ex}

\begin{ex} %[5D10C61]
Một vận động viên leo núi dự định leo lên đỉnh núi cao $1\,200\text{ m}$. Vận tốc khi leo lên là $200\text{ m/giờ}$, vận tốc khi xuống núi là $400\text{ m/giờ}$. Người đó nghỉ trên đỉnh núi $1$ giờ $30$ phút. Tổng thời gian của cả chuyến leo núi là
\choice
{$9$ giờ}
{$10$ giờ}
{\True $10$ giờ $30$ phút}
{$11$ giờ}
\loigiai{
Thời gian leo lên: $1\,200 : 200 = 6$ giờ.\\
Thời gian leo xuống: $1\,200 : 400 = 3$ giờ.\\
Thời gian nghỉ: $1$ giờ $30$ phút.\\
Tổng thời gian: $6 + 3 + 1$ giờ $30$ phút $= 10$ giờ $30$ phút.\\
\textbf{Đáp án của câu hỏi là C}
}
\end{ex}

\textbf{Phần 2. Trắc nghiệm đúng sai}
\setcounterpageref{ex}{0}

\begin{ex} %[5D10H61]
Xét tính đúng sai về ước lượng vận tốc trong thực tế:
\choiceTF
{\True Vận tốc của máy bay chở khách thông thường khoảng $800 - 900\text{ km/h}$.}
{Vận tốc người chạy bộ thông thường có thể đạt $100\text{ km/h}$.}
{\True Một người đi xe đạp điện với vận tốc $25\text{ km/h}$ đi $5\text{ km}$ hết $12$ phút.}
{\True Vận tốc âm thanh trong không khí xấp xỉ $340\text{ m/s}$.}
\loigiai{
\begin{itemize}
    \item Máy bay chở khách: 800 - 900 km/h. Đúng.
    \item Người chạy bộ không thể đạt 100 km/h (tối đa chỉ khoảng 35-40 km/h đối với VĐV thế giới). Sai.
    \item $5 : 25 = 0{,}2\text{ giờ} = 12$ phút. Đúng.
    \item Vận tốc âm thanh: ~340 m/s. Đúng.
\end{itemize}
\begin{itemize}
    \item a) \textbf{Đúng}.
    \item b) \textbf{Sai}.
    \item c) \textbf{Đúng}.
    \item d) \textbf{Đúng}.
\end{itemize}
}
\end{ex}

\begin{ex} %[5D10V61]
Một gia đình lái xe ô tô đi nghỉ mát từ tỉnh A đến tỉnh B cách nhau $240\text{ km}$. Họ dự định lái xe với vận tốc $60\text{ km/h}$ và nghỉ dọc đường $45$ phút. Xét tính đúng sai:
\choiceTF
{\True Thời gian lái xe thực tế là $4$ giờ.}
{\True Tổng thời gian cả đi và nghỉ là $4$ giờ $45$ phút.}
{Nếu xuất phát lúc $6$ giờ $30$ phút thì đến nơi lúc $11$ giờ.}
{\True Nếu xuất phát lúc $6$ giờ $30$ phút thì đến nơi lúc $11$ giờ $15$ phút.}
\loigiai{
\begin{itemize}
    \item Thời gian lái xe: $240 : 60 = 4$ giờ. Đúng.
    \item Tổng thời gian: $4$h $+ 45$p $= 4$h $45$p. Đúng.
    \item Đến nơi: $6$h $30$p $+ 4$h $45$p $= 11$h $15$p $\neq 11$h. Ý 3 Sai, Ý 4 Đúng.
\end{itemize}
\begin{itemize}
    \item a) \textbf{Đúng}.
    \item b) \textbf{Đúng}.
    \item c) \textbf{Sai}.
    \item d) \textbf{Đúng}.
\end{itemize}
}
\end{ex}

\textbf{Phần 3. Trả lời ngắn}
\setcounterpageref{ex}{0}

\begin{ex} %[5D10H61]
Nhà Minh cách trường $2{,}4\text{ km}$. Minh đi xe đạp hết $12$ phút. Tính vận tốc của Minh theo $\text{km/h}$.
\shortans{12}
\loigiai{
\textbf{Đáp số:} $12$\\
$12$ phút $= 0{,}2$ giờ $\Rightarrow 2{,}4 : 0{,}2 = 12\text{ km/h}$.
}
\end{ex}

\begin{ex} %[5D10H61]
Một người đi xe máy trong $45$ phút với vận tốc $40\text{ km/h}$. Tính quãng đường người đó đi được theo ki-lô-mét.
\shortans{30}
\loigiai{
\textbf{Đáp số:} $30$\\
$40 \times 0{,}75 = 30\text{ km}$.
}
\end{ex}

\begin{ex} %[5D10V61]
Một vận động viên bơi lội bơi cự li $200\text{ m}$ hết $2$ phút $30$ giây. Tính vận tốc bơi theo mét trên giây (dùng dấu phẩy nếu có).
\shortans{1,33 hoặc \dfrac{4}{3}}
\loigiai{
\textbf{Đáp số:} $1{,}33$ hoặc $\dfrac{4}{3}$\\
$2$ phút $30$ giây $= 150$ giây $\Rightarrow 200 : 150 = 1{,}33\text{ m/s}$.
}
\end{ex}

\begin{ex} %[5D10V61]
Một người đi xe đạp với vận tốc $15\text{ km/h}$. Người đó muốn đi quãng đường $22{,}5\text{ km}$ thì mất bao nhiêu phút?
\shortans{90}
\loigiai{
\textbf{Đáp số:} $90$\\
$22{,}5 : 15 = 1{,}5\text{ giờ} = 90\text{ phút}$.
}
\end{ex}

\begin{ex} %[5D10V61]
Một tàu thủy chạy trên biển với vận tốc $18$ hải lí/giờ ($1$ hải lí $\approx 1{,}852\text{ km}$). Tính quãng đường tàu chạy trong $3$ giờ theo ki-lô-mét.
\shortans{100,008}
\loigiai{
\textbf{Đáp số:} $100{,}008$\\
$18 \times 3 \times 1{,}852 = 54 \times 1{,}852 = 100{,}008\text{ km}$.
}
\end{ex}

\begin{ex} %[5D10C61]
Một nhóm học sinh đi dã ngoại. Chặng 1 đi xe đạp dài $12\text{ km}$ với vận tốc $12\text{ km/h}$. Chặng 2 đi bộ dài $4\text{ km}$ với vận tốc $4\text{ km/h}$. Nghỉ giữa hai chặng $30$ phút. Tính tổng thời gian chuyến đi dã ngoại theo phút.
\shortans{150}
\loigiai{
\textbf{Đáp số:} $150$\\
Chặng 1: $1$ giờ $= 60$ phút.\\
Chặng 2: $1$ giờ $= 60$ phút.\\
Nghỉ: $30$ phút.\\
Tổng: $60 + 60 + 30 = 150$ phút.
}
\end{ex}
